% Motivation and learning outcomes.
% The outcomes are the outcomes of Day 11 in the syllabus.
\begin{frame}{Why this day matters}
  \begin{itemize}
    \item Days 1 to 10 ran a model on one device. Week 3 connects the devices:
          a camera node sends images, a second node runs the model, and a
          third node shows the results.
    \item A network adds delay, loses packets, and has a limited bandwidth. A
          video stream at 30 frames each second can need more bandwidth than
          the Wi-Fi of a classroom gives to one device.
    \item A detector on a network stream sees an image that is old. The age of
          the image comes from the camera, the encoder, the network, the
          buffers, and the model. The latency of the model is only one part.
    \item In the lab you find the devices on the lab network, publish the
          camera of the Raspberry Pi as an RTSP stream, run the Day 8 detector
          on the stream, and measure the end-to-end latency. Then you stream
          the camera of the XIAO and compare.
  \end{itemize}
\end{frame}

\begin{frame}{Learning outcomes}
  After this day you can:
  \begin{itemize}
    \item Explain when to use TCP and when to use UDP.
    \item Describe an RTSP session and the role of RTP and H.264.
    \item Set up an RTSP server on the Raspberry Pi.
    \item Read a video stream in Python and run inference on it.
    \item Measure the end-to-end latency of a stream.
  \end{itemize}
\end{frame}
